\fichatitulo{14 · CONTEXTO}{Artigo · 2016}{Informação e engajamento parlamentar}
{\small Bernardes; Leston-Bandeira. \textit{Information vs Engagement in Parliamentary Websites -- A Case Study of Brazil and the UK}. Artigo · 2016. \href{https://revistas.ufpr.br/rsp/article/download/48709/29300}{Fonte consultada}.\par}

\campo{Problema e objetivo}
Como os portais parlamentares articulam informação e participação?

\campo{Principal contribuição · síntese interpretativa}
Distingue acesso à informação institucional, informação sobre atividades e possibilidades de engajamento, comparando Brasil e Reino Unido.

\campo{Método / natureza da evidência}
Estudo comparativo das câmaras baixas, com análise de portais e entrevistas com funcionários.

\campo{Resultado ou argumento principal}
Predomina a oferta de informação; os mecanismos de interação não demonstram, por si, participação efetiva no processo legislativo.

\campo{Excertos-chave · seleção literal}
\begin{itemize}
\excerto{1}{there is still little evidence of enabling citizen participation in the legislative process}{Resumo.}
\excerto{2}{information about parliamentary activity}{Resumo.}
\end{itemize}

\campo{Limitações e análise crítica}
Análise da ficha: o estudo descreve instituições e tecnologias de seu período. Não avalia RAG nem mede o efeito causal de uma interface sobre a participação.

\campo{Aproveitamento na dissertação · proposta}
Introdução: justificar o problema de acesso sem equiparar desempenho técnico a democratização.

\registro{Fonte consultada nesta preparação: Resumo, seção III e conclusão (VI). Chave: \texttt{bandeiraBernardes2016information}.}
